\section{Splunk and SPL Fundamentals}

Splunk is one of the most widely used platforms for log analysis, threat hunting, and detection engineering. Its strength lies in its ability to ingest large volumes of telemetry from diverse sources and make that data searchable in near real time. For a CTI analyst or threat hunter, Splunk becomes the primary interface for understanding system behavior, identifying anomalies, and validating hypotheses. This section introduces the core concepts of Splunk and its query language, SPL, in a way that supports long-term analytical work.

At its core, SPL (Search Processing Language) is designed to filter, transform, and analyze event data. Most SPL searches follow a simple pattern: start with an index, apply filters, and then use commands such as \texttt{stats}, \texttt{table}, or \texttt{eval} to shape the results. The goal is not to memorize every command, but to understand how to combine them to answer investigative questions. Because many organizations also rely on Microsoft Sentinel, each SPL example here is paired with an equivalent KQL version to reinforce cross-platform skills.

\subsection*{Basic Search Structure}

Before diving into threat hunting examples, it is useful to understand the basic structure of an SPL search. The following example retrieves Windows process creation events from Sysmon:

\textbf{SPL Example}
\begin{lstlisting}
index=sysmon EventCode=1
| table _time, Computer, User, Image, CommandLine
\end{lstlisting}

\textbf{KQL Equivalent}
\begin{lstlisting}
Sysmon
| where EventID == 1
| project TimeGenerated, Computer, User, Image, CommandLine
\end{lstlisting}

This simple pattern—filter, then shape—is the foundation of nearly every SPL and KQL query you will write.

\subsection*{Filtering and Field Matching}

Filtering is one of the most important skills in SPL. Analysts often begin with broad searches and gradually narrow them down as they learn more about the activity. The following example searches for PowerShell executions that include encoded commands, a common sign of obfuscation:

\textbf{SPL Example}
\begin{lstlisting}
index=wineventlog EventCode=4104
| search CommandLine="*EncodedCommand*"
\end{lstlisting}

\textbf{KQL Equivalent}
\begin{lstlisting}
PowerShellEvent
| where EventID == 4104
| where CommandLine contains "EncodedCommand"
\end{lstlisting}

These types of queries are frequently used in threat hunting, detection engineering, and CTI validation.

\subsection*{Aggregation and Pattern Analysis}

Threat hunters often need to identify unusual patterns rather than single events. Aggregation commands such as \texttt{stats} allow analysts to group events and look for anomalies. The following example identifies accounts with repeated failed logons:

\textbf{SPL Example}
\begin{lstlisting}
index=wineventlog EventCode=4625
| stats count by AccountName, WorkstationName
| where count > 10
\end{lstlisting}

\textbf{KQL Equivalent}
\begin{lstlisting}
SecurityEvent
| where EventID == 4625
| summarize Count=count() by Account, Workstation
| where Count > 10
\end{lstlisting}

This type of aggregation is foundational for detecting brute-force attempts, credential misuse, and lateral movement.

\subsection*{Process Relationships}

Understanding parent-child process relationships is essential for detecting suspicious execution chains. SPL makes this easy by correlating fields such as \texttt{ParentImage} and \texttt{Image}. The following example highlights unusual parent-child pairs:

\textbf{SPL Example}
\begin{lstlisting}
index=sysmon EventCode=1
| stats count by ParentImage, Image
| where count < 5
\end{lstlisting}

\textbf{KQL Equivalent}
\begin{lstlisting}
Sysmon
| where EventID == 1
| summarize Count=count() by ParentImage, Image
| where Count < 5
\end{lstlisting}

This is a common technique in threat hunting to identify rare or suspicious execution paths.

\subsection*{Why These Fundamentals Matter}

These examples form the foundation of nearly all detection and hunting work. Whether you are validating CTI findings, building dashboards, or investigating an alert, the ability to quickly filter, aggregate, and correlate telemetry is essential. As the handbook progresses, these fundamentals will be expanded into full detection patterns, hunting playbooks, and dashboard designs that reflect real-world analyst workflows.
